%% LyX 2.5.1 created this file.  For more info, see https://www.lyx.org/.
%% Do not edit unless you really know what you are doing.
\documentclass[10pt,english,t]{beamer}
\usepackage{lmodern}
\usepackage[T1]{fontenc}
\usepackage[utf8]{inputenc}
\usepackage{amsbsy}
\usepackage{amstext}
\usepackage{amssymb}
\usepackage{graphicx}
\usepackage[authoryear]{natbib}
\PassOptionsToPackage{normalem}{ulem}
\usepackage{ulem}

\makeatletter
%%%%%%%%%%%%%%%%%%%%%%%%%%%%%% Textclass specific LaTeX commands.
% this default might be overridden by plain title style
\newcommand\makebeamertitle{\frame{\maketitle}}%
% (ERT) argument for the TOC
\AtBeginDocument{%
  \let\origtableofcontents=\tableofcontents
  \def\tableofcontents{\@ifnextchar[{\origtableofcontents}{\gobbletableofcontents}}
  \def\gobbletableofcontents#1{\origtableofcontents}
}

%%%%%%%%%%%%%%%%%%%%%%%%%%%%%% User specified LaTeX commands.
\usepackage{tikz}
\usetikzlibrary{positioning}
\usepackage{appendixnumberbeamer}

\usetheme[progressbar=frametitle,block=fill]{metropolis}

% margin
\setbeamersize{text margin right=1.5cm}

% colors
\colorlet{DarkRed}{red!70!black}
\setbeamercolor{normal text}{fg=black}
\setbeamercolor{alerted text}{fg=DarkRed}
\setbeamercolor{progress bar}{fg=DarkRed}
\setbeamercolor{button}{bg=DarkRed}

% width of seperators
\makeatletter
\setlength{\metropolis@titleseparator@linewidth}{1pt}
\setlength{\metropolis@progressonsectionpage@linewidth}{1pt}
\setlength{\metropolis@progressinheadfoot@linewidth}{1pt}
\makeatother

% new alert block
\newlength\origleftmargini
\setlength\origleftmargini\leftmargini
\setbeamertemplate{itemize/enumerate body begin}{\setlength{\leftmargini}{4mm}}
\let\oldalertblock\alertblock
\let\oldendalertblock\endalertblock
\def\alertblock{\begingroup \setbeamertemplate{itemize/enumerate body begin}{\setlength{\leftmargini}{\origleftmargini}} \oldalertblock}
\def\endalertblock{\oldendalertblock \endgroup}
\setbeamertemplate{mini frame}{}
\setbeamertemplate{mini frame in current section}{}
\setbeamertemplate{mini frame in current subsection}{}
\setbeamercolor{section in head/foot}{fg=normal text.bg, bg=structure.fg}
\setbeamercolor{subsection in head/foot}{fg=normal text.bg, bg=structure.fg}

% footer
\makeatletter
\setlength{\metropolis@frametitle@padding}{1.6ex}
\setbeamertemplate{footline}{%
    \begin{beamercolorbox}[colsep=1.5pt]{upper separation line head}
    \end{beamercolorbox}
    \begin{beamercolorbox}{section in head/foot}
      \vskip1pt\insertsectionnavigationhorizontal{\paperwidth}{}{\hskip0pt plus1filll \insertframenumber{} / \inserttotalframenumber \hskip2pt}\vskip3pt% 
    \end{beamercolorbox}%
    \begin{beamercolorbox}[colsep=1.5pt]{lower separation line head}
    \end{beamercolorbox}
}
\makeatother

% toc
\setbeamertemplate{section in toc}{\hspace*{1em}\inserttocsectionnumber.~\inserttocsection\par}
\setbeamertemplate{subsection in toc}{\hspace*{2em}\inserttocsectionnumber.\inserttocsubsectionnumber.~\inserttocsubsection\par}
% Added by lyx2lyx
\setlength{\parskip}{\smallskipamount}
\setlength{\parindent}{0pt}

\makeatother

\usepackage{babel}
\begin{document}
\title{14. Exam and Perspectives}
\subtitle{Adv. Macro: Heterogenous Agent Models} 
\author{Raphaël Huleux}
\date{2026}

{
\setbeamertemplate{footline}{} 
\begin{frame}
\maketitle
\end{frame}
}

\addtocounter{framenumber}{-1}

\begin{frame}<beamer>
\frametitle{Plan}
\tableofcontents[]
\end{frame}

\section{Learning outcomes}
\begin{frame}{Knowledge, Skills and Competencies}
\begin{itemize}
\item \textbf{What do you need to know?}
\begin{enumerate}
\item Understand and be able to use the \emph{computational techniques}
to solve and simulate heterogeneous agent models
\item Understand and be able to discuss the \emph{economic insights} heterogeneous
agent models can provide
\end{enumerate}
\item \textbf{Next: }Learning outcomes in Knowledge, Skills and Competencies
\end{itemize}
\end{frame}
%
\begin{frame}{Knowledge}
\begin{enumerate}
\item {\small Account for, formulate and interpret precautionary saving models }{\small\par}
\item {\small Account for stochastic and non-stochastic simulation methods }{\small\par}
\item {\small Account for, formulate and interpret general equilibrium models
with ex ante and ex post heterogeneity, idiosyncratic and aggregate
risk, and with and without pricing frictions }{\small\par}
\item {\small Discuss the difference between the stationary equilibrium,
the transition path and the dynamic equilibrium }{\small\par}
\item {\small Discuss the relationship between various equilibrium concepts
and their solution methods}{\small\par}
\item {\small Identify and account for methods for analyzing the dynamic
distributional effects of long-run policy (e.g. taxation and social
security) and short-run policy (e.g. monetary and fiscal policy)}{\small\par}
\end{enumerate}
\end{frame}
%
\begin{frame}{Skills}
\begin{enumerate}
\item {\small Solve precautionary saving problems with dynamic programming
and simulate behavior with stochastic and non-stochastic techniques }{\small\par}
\item {\small Solve general equilibrium models with ex ante and ex post heterogeneity,
idiosyncratic and aggregate risk, and with and without pricing frictions
(stationary equilibrium, transition path, dynamic equilibrium) }{\small\par}
\item {\small Analyze dynamics of income and wealth inequality }{\small\par}
\item {\small Analyze transitional and permanent structural changes (e.g.
inequality trends and the long-run decline in the interest rate) }{\small\par}
\item {\small Analyze the dynamic distributional effects of long-run policy
(e.g. taxation and social security) and short-run policy (e.g. monetary
and fiscal policy)}{\small\par}
\end{enumerate}
\end{frame}
%
\begin{frame}{Competencies}
\begin{enumerate}
\item {\small Independently formulate, discuss and assess research on both
the causes and effects of heterogeneity and risk for both long-run
and short-run outcomes }{\small\par}
\item {\small Discuss and assess the importance of how heterogeneity and
risk is modeled for questions about both long-run and short-run dynamics}{\small\par}
\end{enumerate}
\end{frame}
%

\section{Overview and perspectives}
\begin{frame}{Single-agent problems (>>partial equilibrium<<, PE)}
\begin{itemize}
\item \textbf{Focus: }Consumption-saving (and labor supply) of households
\item \textbf{Solution methods:}
\begin{enumerate}
\item Value-Function Iteration (VFI)
\item Endogenous Gridpoint method (EGM)
\end{enumerate}
\item \textbf{Simulation methods:}
\begin{enumerate}
\item Monte Carlo
\item Histogram
\end{enumerate}
\item \textbf{Perspectives:}
\begin{enumerate}
\item Multiple assets (liquidity, riskiness, durables, housing)
\item Similar analysis of \emph{firm problems}
\item Closer link to \emph{empirical observations} (life-cycle dynamics,
informational frictions, behavioral biases, social preferences, norms)
\item More \emph{complex models} (solved with e.g. >>deep learning<<)
\end{enumerate}
\end{itemize}
\end{frame}
%
\begin{frame}{Stationary equilibrium}
\begin{itemize}
\item \textbf{Stationary:} \emph{Time-constant aggregate variables and distribution}
\begin{enumerate}
\item \textbf{Fixed point problem:} Choose all parameters (and normalizations)
and search for prices or aggregate quantities
\item \textbf{Combined with calibration: }Some parameters set residually
and/or there are targets beyond market clearing (e.g. MPC)
\end{enumerate}
\item \textbf{Comparative statics: }How does equilibrium outcomes change
when one or more parameters change?
\begin{enumerate}
\item Keep all other parameters (and normalizations) fixed
\item Re-solve fixed point problem for prices or aggregate quantities
\end{enumerate}
\item \textbf{Perspective:}
\begin{enumerate}
\item Always a good starting point for analyzing dynamics
\item Understanding of average outcomes in >>stable periods<<
\item Deeper >>welfare analysis<< (maximization of social welfare function)
\end{enumerate}
\end{itemize}
\end{frame}
%
\begin{frame}{Dynamic equilibrium: Aggregate risk}
\begin{itemize}
\item \textbf{Canonical state-space form with aggregate risk:}
\begin{enumerate}
\item States summarize all relevant information
\item Policy functions (idiosyncratic + aggregate) are functions of states{\small
\[
\left[\begin{array}{c}
\boldsymbol{X}_{t+1},\boldsymbol{Z}_{t+1},\underline{\boldsymbol{D}}_{t+1},\boldsymbol{Y}_{t}\end{array}\right]^{\prime}=\mathcal{M}\left(\left[\begin{array}{c}
\boldsymbol{X}_{t},\boldsymbol{Z}_{t},\underline{\boldsymbol{D}}_{t}\end{array}\right]^{\prime},\boldsymbol{\epsilon}_{t}\right)
\]
}{\small\par}

Expectations are hard:
\begin{itemize}
\item Depend on current state (high-dimensional)
\item Requires knowing $\mathcal{M}$ (perceived law-of-motion, PLM)
\end{itemize}
\end{enumerate}
\item \textbf{Fluctuations:}
\begin{enumerate}
\item Volatility of shocks $\rightarrow0$: Deterministic steady state (DSS)
\item Realization of shocks $\rightarrow0$: Stochastic steady state (SSS)
\item Otherwise: Convergence to \emph{ergodic distribution}
\end{enumerate}
\item \textbf{Moving average form: }History of aggregate shocks (without
any truncation) determine current states $\Rightarrow$
\[
\left[\begin{array}{c}
\boldsymbol{X}_{t+1},\boldsymbol{Z}_{t+1},\underline{\boldsymbol{D}}_{t+1},\boldsymbol{Y}_{t}\end{array}\right]^{\prime}=\tilde{\mathcal{M}}\left(\boldsymbol{\epsilon}_{t},\boldsymbol{\epsilon}_{t-1},\boldsymbol{\epsilon}_{t-2},\dots\right)
\]

\textbf{Solve on MA form:} Azinovic-Yand and Žemlička (2025, Sep 16)
\end{itemize}
\end{frame}
%
\begin{frame}{Sequence space}
\begin{itemize}
\item \textbf{Linearization in state space:}{\small
\[
\left[\begin{array}{c}
\boldsymbol{X}_{t+1},\boldsymbol{Z}_{t+1},\underline{\boldsymbol{D}}_{t+1},\boldsymbol{Y}_{t}\end{array}\right]^{\prime}=\boldsymbol{A}\left[\begin{array}{c}
\boldsymbol{X}_{t},\boldsymbol{Z}_{t},\underline{\boldsymbol{D}}_{t}\end{array}\right]^{\prime}+\boldsymbol{B}\boldsymbol{\epsilon}_{t}
\]
}{\small\par}
\item \textbf{Sequence space: }Perfect foresight wrt. aggregate variables

Behavior depend on:
\begin{enumerate}
\item Idiosyncratic state variables
\item Current value and expected distribution of relevant >>prices<<\\
$\Rightarrow$ collapses to sequence (time-path) of relevant >>prices<<
\end{enumerate}
\item \textbf{MIT shock: }One-time shock in a perfect foresight model
\item \textbf{Insight: }\emph{The first-order IRF from an MIT shock is }\emph{\uline{equivalent}}\emph{
to the IRF in a model with aggregate risk, which is linearized in
the aggregate variables} (Boppart et. al., 2018)
\end{itemize}
\end{frame}
%
\begin{frame}{GE and Household Jacobians}
\begin{itemize}
\item \textbf{Target equation system }with implicit truncation
\[
\boldsymbol{H}(\boldsymbol{U},\boldsymbol{Z})=\boldsymbol{0}
\]
\item \textbf{First-order IRF to $d\boldsymbol{Z}$: }
\begin{align*}
\boldsymbol{H}(\boldsymbol{U},\boldsymbol{Z})=0\Rightarrow\boldsymbol{H}_{\boldsymbol{U}}d\boldsymbol{U}+\boldsymbol{H}_{\boldsymbol{Z}}d\boldsymbol{Z}=0 & \Leftrightarrow d\boldsymbol{U}=-\boldsymbol{H}^{-1}_{\boldsymbol{U}}\boldsymbol{H}_{\boldsymbol{Z}}d\boldsymbol{Z}
\end{align*}
All other variables: Easy with $d\boldsymbol{U}$ (and $d\boldsymbol{Z}$)
\item \textbf{Non-linear transition path: }Fully solve $\boldsymbol{H}(\boldsymbol{U},\boldsymbol{Z}_{ss}+d\boldsymbol{Z})=\boldsymbol{0}$
with quasi-Newton method initially using $\boldsymbol{H}_{\boldsymbol{U}}$
for updating guesses
\item \textbf{Calculation of derivatives:}
\begin{enumerate}
\item Derivatives of aggregate equations found with numerical derivative
\item Derivatives of household block found with fake-news algorithm
\begin{enumerate}
\item Backward: Only time to shock matters
\item Forward: Apply chain-rule to reduce computation burden\\
(parallelization is an alternative, especially in more complex models)
\end{enumerate}
\end{enumerate}
\end{itemize}
\end{frame}
%
\begin{frame}{Aggregate consumption function}
\begin{itemize}
\item \textbf{Household problem}
\begin{align*}
v_{t}(z_{t},a_{t-1}) & =\max_{c_{t}}\frac{c^{1-\sigma}_{t}}{1-\sigma}+\beta\mathbb{E}_{t}\left[v_{t+1}(z_{t+1},a_{t})\right]\\
\text{s.t. }a_{t}+c_{t} & =(1+r_{t})a_{t-1}+(1-\tau_{t})w_{t}\ell_{t}z_{t}\\
\log z_{t+1} & =\rho_{z}\log z_{t}+\psi_{t+1}\,\,\,,\psi_{t}\sim\mathcal{N}(\mu_{\psi},\sigma_{\psi}),\,\mathbb{E}[z_{t}]=1\\
a_{t} & \geq0
\end{align*}
\item \textbf{Consumption function:
\begin{align*}
C^{hh}_{t} & =C^{hh}\left(\{r_{s},\tau_{s},w_{s},\ell_{s}\}_{s\geq0}\right)\Rightarrow\\
\boldsymbol{C}^{hh} & =C^{hh}(\boldsymbol{r},\boldsymbol{\tau},\boldsymbol{w},\boldsymbol{\ell})\Rightarrow\\
d\boldsymbol{C}^{hh} & =\boldsymbol{M}^{r}d\boldsymbol{r}+\boldsymbol{M}^{\tau}d\boldsymbol{\tau}+\boldsymbol{M}^{w}d\boldsymbol{w}+\boldsymbol{M}^{\ell}d\boldsymbol{\ell}
\end{align*}
}where $\boldsymbol{M}^{x}=\left[\partial C_{t}/\partial x_{s}\right]_{ts}$
or with $Z_{t}\equiv(1-\tau_{t})w_{t}\ell_{t}$
\[
d\boldsymbol{C}^{hh}=\boldsymbol{M}^{r}d\boldsymbol{r}+\boldsymbol{M}d\boldsymbol{Z}
\]
\end{itemize}
\end{frame}
%
\begin{frame}{Intertemporal Keynesian Cross I}
\begin{itemize}
\item \textbf{Linear production:} $Y_{t}=w_{t}L_{t}$ with $\int\ell_{t}d\boldsymbol{D}_{t}$
\item \textbf{Taxes:} $T_{t}=\tau_{t}w_{t}L_{t}$
\item \textbf{Unions:} $\ell_{t}=L_{t}$
\item \textbf{Disposable income:} $Z_{t}=Y_{t}-T_{t}$
\item \textbf{Market clearing:
\begin{align*}
d\boldsymbol{Y} & =d\boldsymbol{G}+d\boldsymbol{C}^{hh}\\
 & =d\boldsymbol{G}+\boldsymbol{M}^{r}d\boldsymbol{r}+\boldsymbol{M}d\boldsymbol{Z}\\
 & =d\boldsymbol{G}+\boldsymbol{M}^{r}d\boldsymbol{r}+\boldsymbol{M}\left[d\boldsymbol{Y}-d\boldsymbol{T}\right]
\end{align*}
}Requirement: $\sum^{\infty}_{t=0}(1+r_{ss})^{-t}(dG_{t}-dT_{t})=0$

Alternative: Specify explicit policy rule for $d\boldsymbol{T}$
\end{itemize}
\end{frame}
%
\begin{frame}{Intertemporal Keynesian Cross II}
\begin{itemize}
\item \textbf{Market clearing:
\begin{align*}
d\boldsymbol{Y} & =d\boldsymbol{G}+\boldsymbol{M}^{r}d\boldsymbol{r}+\boldsymbol{M}\left[d\boldsymbol{Y}-d\boldsymbol{T}\right]
\end{align*}
}
\item \textbf{Intertemporal Keynesian Cross: }$d\boldsymbol{r}=0$ implies\textbf{
\begin{align*}
d\boldsymbol{Y} & =d\boldsymbol{G}+\boldsymbol{M}\left[d\boldsymbol{Y}-d\boldsymbol{T}\right]\Leftrightarrow\\
\left[\boldsymbol{I}-\boldsymbol{M}\right]d\boldsymbol{Y} & =d\boldsymbol{G}-\boldsymbol{M}d\boldsymbol{T}\Leftrightarrow\\
d\boldsymbol{Y} & =\mathcal{M}\left[d\boldsymbol{G}-\boldsymbol{M}d\boldsymbol{T}\right]
\end{align*}
}

where $\mathcal{M}\left[\boldsymbol{I}-\boldsymbol{M}\right]=\boldsymbol{I}$
(left-inverse)
\begin{itemize}
\item \textbf{Uniqueness:} Can be proven (single $\mathcal{M}$ imply bounded
solution)
\item \textbf{Takeaway:} iMPCs in $\boldsymbol{M}$ (+ policy rule determining
$d\boldsymbol{T}$) are sufficient statistics for fiscal multipliers.
\end{itemize}
\end{itemize}
\end{frame}
%
\begin{frame}{From HANC to HANK}
\begin{itemize}
\item \textbf{HANC: }\emph{The real interest must ensure market clearing}
\begin{enumerate}
\item Production: Endogenous labor supply and/or capital accumulation
\item Exchange: Endowment or exogenous labor supply $\div$ capital acc.
\end{enumerate}
\item \textbf{Canonical HANK with demand-determined labor supply:}\emph{}\\
\emph{CB controls real rate and output adjust to market clearing}
\begin{enumerate}
\item Households and government: Exogenous labor supply without capital
accumulation $\sim$ HANC exchange economy 
\item Central bank and production: Taylor-rule and New Keynesian Wage Phillips
Curve link real rate to output
\item Labor market and union/search-and-matching: Link output to idiosyncratic
income process
\end{enumerate}
\item \textbf{Sticky prices, but fully flexible wages}: Empirically unrealistic
large wealth effects on labor supply
\item \textbf{More sticky wages than profits:} Pro-cyclical profits
\end{itemize}
\end{frame}
%
\begin{frame}{Transition path}
\begin{itemize}
\item \textbf{Transition path}
\begin{enumerate}
\item Temporary shock: Non-linear IRF
\item Permanent shock: Transition to new stationary equilibrium
\end{enumerate}
\item \textbf{Welfare analysis conditional on initial state}

Example: Looser borrowing constraint
\begin{enumerate}
\item Larger choice set $\Rightarrow$ higher utility
\item Lower utility in stationary equilibrium possible due to dissaving... 
\end{enumerate}
\textbf{Details:} Policy tools (e.g. restricted by functional forms)
and degree of commitment (time-consistency)
\item \textbf{More complex models: }New mechanisms? Inequality outcomes?
\begin{enumerate}
\item Life-cycle and demographics
\item Spatial dimension
\item Structural transformation {\footnotesize (agriculture $\rightarrow$
manufacturing $\rightarrow$ services)}{\footnotesize\par}
\item AI and automation
\item Climate
\end{enumerate}
\item \textbf{Computationally: }\emph{Imprecise solution become an issue}
\end{itemize}
\end{frame}
%
\begin{frame}{Limitation of sequence space approach}
\begin{itemize}
\item \textbf{First-order approximation:}
\begin{enumerate}
\item Around deterministic steady state (DSS)
\item No state and size-dependence of IRFs

(the non-linear transition path to a MIT shock has size-dependence)
\end{enumerate}
\item \textbf{Second-order approximation:}
\begin{enumerate}
\item Around stochastic steady state (SSS)
\item Can be derived from MIT shock {\footnotesize (Auclert, Rognlie, Straub,
SED2025)}{\footnotesize\par}
\item Approximation need not be better
\end{enumerate}
\item \textbf{Full dynamic equilibrium:}
\begin{enumerate}
\item Distribution affected by previous aggregate shocks\\
(e.g. distribution of fixed rate mortgage contracts)
\item Feedback effects from policy\\
(e.g. lower self-insurance with counter-cyclical UI)
\item Directly allow for fully non-linear dynamics\\
(e.g. zero lower bound, capacity constraints etc.)
\end{enumerate}
\end{itemize}
\end{frame}
%
\begin{frame}{Business cycle questions}
\begin{itemize}
\item \textbf{Importance of heterogeneity for business cycles:}
\begin{enumerate}
\item Some aggregate neutrality results - still distributional concerns
\item Size of mechanisms are different - cash-flow effects important
\item High MPC and precautionary saving become of central importance
\end{enumerate}
\item \textbf{Investment: }Still not strongly demand-driven
\begin{enumerate}
\item Positive correlation from investment shocks\\
(due to cash-flow effects and less precautionary saving)
\item Negative correlation from consumption (due higher real rate)
\end{enumerate}
\item \textbf{Other open questions:}
\begin{enumerate}
\item Imperfect information and behavioral biases {\footnotesize (e.g. wage
negotiation in inflation surges, experimental/survey evidence)}{\footnotesize\par}
\item Asymmetric fluctuations {\footnotesize (plucking models, efficient
unemployment) }{\footnotesize\par}
\item Fundamental uncertainty {\footnotesize (inflation due to demand or
supply?)}{\footnotesize\par}
\item Coordination problems {\footnotesize (Walras $\rightarrow$ search-and-matching
$\rightarrow$ ? $\sim$ agent-based)}{\footnotesize\par}
\end{enumerate}
\end{itemize}
\end{frame}
%

\section{Exam}
\begin{frame}{Exam}
\begin{itemize}
\item <+->\textbf{Exam: }Similar to assignment I-II
\begin{enumerate}
\item \textbf{Starting point: }
\begin{enumerate}
\item Model definition in .py-files which can run
\item Notebook with pre-defined structure and some plotting functions
\end{enumerate}
\item \textbf{Code questions:} Add a few lines of code in the right place
\item \textbf{Math questions:} Derive analytical expressions
\item \textbf{Economic question:} Interpret analytical and quantitative
results
\end{enumerate}
\item <+->\textbf{Grading: }In relation to learning outcomes
\begin{enumerate}
\item Correct and well-organized mathematical derivations and code
\item Clear presentation of results (e.g. in figures and tables)
\item Clear and precise text stating the economic conclusions
\end{enumerate}
\item <+->\textbf{Portfolio: }Assignment I-III test the same learning outcomes
\\
$\sim$ floor under not having understood some learning outcome
\item <+->\textbf{Time pressure (36 hours): }More questions than in previous
years
\begin{enumerate}
\item More comprehensive starting point
\item More step-by-step questions
\end{enumerate}
\item <+->\textbf{Code not working: }Describe in words what you have tried
to do
\end{itemize}
\end{frame}
%
\begin{frame}{Hand-in}

\includegraphics[width=1\textwidth]{figs/handin}
\end{frame}
%
\begin{frame}{Preparing for the exam}
\begin{enumerate}
\item Read through slides and study the accompanying code and go through
the exercises
\item Read the documentation for GEModelTools
\item Glance at the central papers
\item Optimize your Assignments I-III
\item Look at the Exam for 2022, 2023 and 2024 (in the Old folder)
\end{enumerate}
\end{frame}
%

\section{GEModelTools}
\begin{frame}{GEModelTools}
\begin{itemize}
\item \textbf{Files: }MODELNAME.py importing blocks.py, household\_problem.py,
steady\_state.py
\item \textbf{Steady state:} .find\_ss() \textbf{$\Rightarrow$ }.ss 
\item \textbf{Tests: }.test\_ss(), .test\_path(),
\item \textbf{Jacobians: }.compute\_jacs(skip\_hh=False,skip\_shocks=False)
\textbf{$\Rightarrow$} 

.jac and .jac\_hh
\item \textbf{First-order IRFs: }.find\_IRFs(shocks=shocks) \textbf{$\Rightarrow$
}.IRF 
\begin{enumerate}
\item shocks = {[}'x'{]} set x to AR(1) with jump\_x and rho\_x
\item shocks = \{dx:dxpath\} set x to ss.x+dxpath
\end{enumerate}
\item \textbf{Non-linear response to MIT shock:} .find\_transition\_path(shocks=shocks)
\textbf{$\Rightarrow$ }.path
\item \textbf{Decomposition of household response:} .decompose\_hh\_path
\item \textbf{Troubleshooting in GEModelTools:} Check examples in notebook
repository + look in documentation or source code
\end{itemize}
\end{frame}
%
\begin{frame}{HANK-sticky-wages\textbackslash simplified}

\emph{Let us look at the code.}
\begin{enumerate}
\item Compute the linear and non-linear IRFs to a AR(1) real rate shock
\item Verify clearing of asset or goods market does not matter
\item Plot household Jacobians
\item Plot linear and non-linear IRFs manually
\item Introduce multiple $\beta$ types
\item Calibrate $B_{ss}$ to an MPC target
\item Change to a balanced budget
\item Compute the linear and non-linear IRF to a custom real rate shock
\item Target a specific output response with government spending
\end{enumerate}
\end{frame}
%

\section{Exam from 2023}
\begin{frame}{Exam from 2023}

\emph{Look at the exam from 2023}
\end{frame}
%

\end{document}
